\section{Curvature pairs and positive lower mass}
\label{sec:elliptic-curvature-endpoints}
\index{Elliptic fibration!positive lower mass}

On a threefold, the lower mass can be bounded below by a curvature condition. Consider a principal elliptic bundle over a K\"ahler surface, with a connection form $\eta$ whose curvature has real and imaginary parts $\alpha$ and $\beta$. If $\alpha^2+\beta^2$ is semipositive, the split metric $g$ built from the base metric and $\eta$ satisfies $\ddc(g^2)=0$ and $\ddc g\leq0$. A cubic energy identity for bounded potentials then shows that no bounded potential has mass below $\int_Xg^3$, and the determinant estimate of \cref{thm:gaud-determinant-endpoint} transfers this lower bound to every Hermitian form on the manifold. The same curvature sign excludes pluriclosed metrics on these manifolds.

In this section $X$ is a connected compact complex threefold and $\mathrm dV$ is a smooth positive probability volume form. Products of bounded potentials are those of \cref{prop:bounded-products}. For a Hermitian form $h$, $\mathscr G_h$ denotes the compact convex set of normalized Gauduchon currents of \cref{sec:gauduchon-linear-duality}, and $k_T$ is the determinant density of $T\in\mathscr G_h$ relative to $\mathrm dV$, as in \cref{thm:gaud-determinant-subcritical}.

\subsection{The bounded cubic identity}

When the square of a Hermitian metric $g$ is $\ddc$-closed, the mass of a bounded $g$-plurisubharmonic function differs from $\int_Xg^3$ by a term that is linear in $\ddc g$. Its coefficient is a positive current playing the role of $\mathrm du\wedge\dc u$.

\begin{lemma}\label{lem:bounded-cubic-lower-mass}
Let $g$ be a Hermitian form with $\ddc(g^2)=0$, and let $u\in\PSH(X,g)$ be bounded. The current
\[
 E_u=\tfrac12\ddc(u^2)-u\ddc u
\]
is positive, and
\begin{equation}\label{eq:bounded-cubic-energy}
 \int_X(g+\ddc u)^3
 =\int_Xg^3-3\int_XE_u\wedge\ddc g.
\end{equation}
If $-\ddc g\geq0$, then $v_-(g)=\int_Xg^3$.
\end{lemma}
\begin{proof}
Locally write $u=p-a$, with $p$ bounded plurisubharmonic and $a$ smooth. Expanding in $p$ and $a$ and using multilinearity, we obtain closed currents $C_r=(\ddc u)^r$, not necessarily positive, with $C_0=1$. The recursive definition of bounded products, expanded in the same way, gives
\[
 \ddc(uC_r)=C_{r+1},\qquad r=0,1,2.
\]
By \cref{lem:weighted-decreasing-products} and invariance under smooth corrections, the weighted product $u\ddc u$ does not depend on the local decomposition. On a smaller ball choose smooth bounded plurisubharmonic $p_j$ decreasing to $p$, and put $u_j=p_j-a$. The same lemma gives $u_j\ddc u_j\to u\ddc u$, while boundedness gives $u_j^2\to u^2$ in $L^1$. Hence
\[
 \mathrm du_j\wedge\dc u_j
 =\tfrac12\ddc(u_j^2)-u_j\ddc u_j\longrightarrow E_u.
\]
The smooth currents on the left are positive, so their weak limit $E_u$ is positive and has order zero.

Now expand the cube. Repeated use of Stokes' theorem for currents gives
\[
 \int_XC_3=0,\qquad
 \int_Xg^2\wedge C_1=\int_Xu\ddc(g^2)=0,
 \qquad
 \int_Xg\wedge C_2=\int_XuC_1\wedge\ddc g.
\]
Since $\ddc g$ is closed, $\int_X\ddc(u^2)\wedge\ddc g=0$, so the last integral equals $-\int_XE_u\wedge\ddc g$. This proves \eqref{eq:bounded-cubic-energy}. If $-\ddc g\geq0$, the correction term is nonnegative, and $u=0$ attains the lower bound $\int_Xg^3$.
\end{proof}

\subsection{The curvature criterion}

We now describe threefolds on which the hypotheses of \cref{lem:bounded-cubic-lower-mass} hold. Let $\pi\colon X\to B$ be a holomorphic principal elliptic bundle over a compact connected K\"ahler surface $(B,\omega)$, with fiber the elliptic curve $E=\C/\Lambda$ and fiber coordinate $z$. Assume that there is a smooth $(1,0)$-form $\eta$ on $X$, invariant under the action of $E$ and restricting to $\mathrm dz$ on every fiber, such that
\[
 \mathrm d\eta=\pi^*F,\qquad F=\alpha+\mathrm i\beta,
\]
where $\alpha,\beta$ are smooth closed real $(1,1)$-forms on $B$. Put
\[
 Q=\alpha^2+\beta^2,\qquad e=\mathrm i\eta\wedge\bar\eta,
 \qquad g=\pi^*\omega+s e,\quad s>0,
 \qquad A_E=\int_E\mathrm i\,\mathrm dz\wedge\mathrm d\bar z.
\]
We call $g$ the split metric.

\begin{theorem}\label{thm:elliptic-curvature-endpoint}
Suppose $Q\geq0$ and $Q\not\equiv0$. The split form $g$ is Hermitian, and its lower mass is
\begin{equation}\label{eq:elliptic-curvature-lower-mass}
 v_-(g)=V=3sA_E\int_B\omega^2.
\end{equation}
Every bounded $u\in\PSH(X,g)$ satisfies
\[
 \int_X(g+\ddc u)^3
 =V+\frac{3s}{2\pi}\int_XE_u\wedge\pi^*Q\geq V.
\]
Every Hermitian form $h\geq c g$, with $c>0$, satisfies \eqref{eq:determinant-background-transport}. Thus the endpoint estimate of \cref{thm:gaud-determinant-endpoint} holds for every Hermitian form on $X$, with a constant depending on the form. The manifold $X$ admits no pluriclosed Hermitian metric.
\end{theorem}
\begin{proof}
We suppress pullbacks in the curvature calculation. Comparing types in $\mathrm d\eta=F$ gives $\partial\eta=0$, $\bar\partial\eta=F$, $\partial\bar\eta=\bar F$ and $\bar\partial\bar\eta=0$. Hence
\[
 \partial e=-\mathrm i\eta\wedge\bar F,\qquad
 \bar\partial e=\mathrm iF\wedge\bar\eta,
 \qquad\ddc e=-\frac1{2\pi}F\wedge\bar F=-\frac1{2\pi}Q.
\]
Since $e^2=0$ and the base has dimension two,
\[
 \ddc g=-\frac{s}{2\pi}\pi^*Q,\qquad
 \ddc(g^2)=-\frac{s}{\pi}\pi^*(\omega\wedge Q)=0,
 \qquad g^3=3s\pi^*(\omega^2)\wedge e.
\]
The form $g$ is positive, because $\pi^*\omega$ is positive on directions transverse to the fibers and $e$ is positive on the vertical line. Fiber integration of $g^3$ gives $V$. \Cref{lem:bounded-cubic-lower-mass} gives the mass identity, and \cref{cor:gaud-determinant-background-comparison}, with $n=3$, gives the endpoint bound. By compactness, every Hermitian form $h$ satisfies $h\geq cg$ for some $c>0$.

Finally, $S=-\ddc g$ is smooth, semipositive and not identically zero. If a Hermitian form $b$ were pluriclosed, Stokes' theorem would give
\[
 0=\int_Xg\wedge\ddc b
   =\int_Xb\wedge\ddc g=-\int_Xb\wedge S<0,
\]
where the last inequality holds because $b$ is strictly positive and $S$ is semipositive and nonzero on an open set.
\end{proof}

\begin{proposition}\label{prop:primitive-curvature-opposite-sign}
The split metric $g$ is balanced if and only if $\omega\wedge\alpha=\omega\wedge\beta=0$. For such a primitive pair one has $Q\leq0$, with strict inequality wherever the pair is nonzero. In particular, a balanced split metric never satisfies the hypothesis $Q\geq0$, $Q\not\equiv0$ of \cref{thm:elliptic-curvature-endpoint}.
\end{proposition}
\begin{proof}
One has $\mathrm d(g^2)=2s\pi^*\omega\wedge\mathrm de$, and $\mathrm de$ has one component containing $\eta$ and one containing $\bar\eta$. These two vertical types are independent, so $\mathrm d(g^2)=0$ if and only if $\omega\wedge F=0$, which is equivalent to the two stated real identities. A primitive real $(1,1)$-form on a surface has eigenvalues $\lambda,-\lambda$ in an $\omega$-unitary frame, so its square is $-\lambda^2\omega^2$. Applying this to $\alpha$ and to $\beta$ gives the sign of $Q$. In this case $\ddc g\geq0$, so \eqref{eq:bounded-cubic-energy} bounds the masses of bounded potentials from above rather than from below, and gives no information on the lower mass.
\end{proof}

\subsection{Negative-line-bundle quotients}
\label{sec:elliptic-line-endpoint-examples}
\index{Elliptic fibration!negative-line-bundle quotient}

The simplest examples of the curvature criterion come from negative line bundles. Dividing the complement of the zero section of a negative line bundle over $\mathbb P^1\times\mathbb P^1$ by a dilation produces a principal elliptic bundle, and the logarithm of the fiber norm gives a connection form whose curvature is a multiple of a K\"ahler form on the base.

Take $B=\mathbb P^1\times\mathbb P^1$, let $\omega_1,\omega_2$ be the pullbacks of the Fubini--Study forms of the two factors, normalized to have integral one, and put
\[
 \begin{gathered}
 \Omega=a\omega_1+b\omega_2,\quad a,b\in\mathbb N_{>0},\qquad
 L=\mathcal O(-a,-b),\\
 X=L^\times/\langle v\mapsto\rho v\rangle,\quad0<\rho<1,
 \end{gathered}
\]
where $L^\times$ is the total space of $L$ with the zero section removed. Equip $L$ with the Hermitian norm induced by the standard norms on the tautological bundles of the two factors. Let $r=\log\|v\|^2$ and set $\eta=\partial r$.

\begin{proposition}\label{prop:elliptic-line-endpoint-example}
The quotient $X$ is a compact complex threefold, and the projection to $B$ is a principal elliptic fibration. For $g=\pi^*\Omega+s\mathrm i\eta\wedge\bar\eta$, $s>0$,
\begin{equation}\label{eq:elliptic-line-lower-mass}
 v_-(g)=24\pi s ab\log(1/\rho).
\end{equation}
Every Hermitian form on $X$ satisfies the endpoint bound of \cref{thm:elliptic-curvature-endpoint} over its full set of normalized Gauduchon currents. The manifold has $H^2(X,\R)\ne0$ and admits no balanced or pluriclosed metric. In particular, it is not diffeomorphic to a Hopf threefold $(\C^3\setminus\{0\})/\langle z\mapsto\lambda z\rangle$, $0<|\lambda|<1$.
\end{proposition}
\begin{proof}
In product coordinates, the squared norm of the standard local frame of $L$ is $(1+|z_1|^2)^a(1+|z_2|^2)^b$. Therefore
\[
 \ddc r=\pi^*\Omega,\qquad \mathrm d\eta=2\pi\mathrm i\pi^*\Omega.
\]
Dilation changes $r$ by a constant, so $\eta$ descends to $X$. The function $r$ changes by the nonzero constant $\log\rho^2$ under the generator, so the action is free and properly discontinuous. The product of the compact unit circle bundle with a fundamental interval for $r$ maps onto $X$, which is therefore compact. Each fiber is $\C^*/\rho^{\mathbb Z}$. Its logarithmic coordinate has periods $\log\rho$ and $2\pi\mathrm i$, and the curvature $F=2\pi\mathrm i\Omega$ has real part $0$ and imaginary part $2\pi\Omega$; hence
\[
 A_E=4\pi\log(1/\rho),\qquad Q=4\pi^2\Omega^2>0,
 \qquad\int_B\Omega^2=2ab.
\]
\Cref{thm:elliptic-curvature-endpoint} now gives \eqref{eq:elliptic-line-lower-mass} and the endpoint bound, and excludes pluriclosed metrics.

To prove $H^2(X,\R)\ne0$, put $A=\eta+\bar\eta$ and $B_1=\mathrm i(\eta-\bar\eta)$. Then $\mathrm dA=0$ and $\mathrm dB_1=-4\pi\pi^*\Omega$. Suppose that $\pi^*(\omega_1-\omega_2)=\mathrm d\zeta$ for a real one-form $\zeta$. Averaging $\zeta$ over the compact group of elliptic translations, we may assume $\zeta$ invariant, so that
\[
 \zeta=\pi^*\xi+(\pi^*f)A+(\pi^*j)B_1
\]
for a one-form $\xi$ and functions $f,j$ on $B$. The components of $\mathrm d\zeta$ with one vertical and one horizontal direction must vanish, which gives $\mathrm df=\mathrm dj=0$. Thus
\[
 \omega_1-\omega_2=\mathrm d\xi-4\pi j\Omega.
\]
Integrating over the two rulings gives $1=-4\pi ja$ and $-1=-4\pi jb$, which is impossible because $a,b>0$. A Hopf threefold as above is diffeomorphic to $S^5\times S^1$, whose second real cohomology vanishes.

Finally, $\pi^*\Omega=-\mathrm dB_1/(4\pi)$ is a nonzero semipositive exact form. Its integral against the closed positive square of a balanced metric would vanish by Stokes' theorem and be strictly positive by positivity. Hence no balanced metric exists.
\end{proof}


\subsection{A Calabi--Eckmann family without LCK metrics}
\label{sec:calabi-eckmann-endpoints}
\index{Calabi--Eckmann manifold!determinant endpoint}

A second construction separates the endpoint criterion from LCK geometry. Let $\tau=A+\mathrm iB_0$, where $A<0$ and $B_0>0$, and let $\C$ act holomorphically on $(\C^2\setminus\{0\})^2$ by
\[
 t\cdot(z,w)=(\mathrm e^t z,\mathrm e^{\tau t}w).
\]
Write $X_\tau$ for the quotient and put $r_1=\log|z|^2$, $r_2=\log|w|^2$. On the projective base let $\omega_1,\omega_2$ be the normalized pulled Fubini--Study forms.

\begin{proposition}\label{prop:calabi-eckmann-endpoint}
The quotient is a compact complex threefold diffeomorphic to $S^3\times S^3$. Projectivization gives a principal elliptic fibration over $\mathbb P^1\times\mathbb P^1$. The form $\eta=\tau\partial r_1-\partial r_2$ descends, and for
\[
 g=\pi^*(\omega_1+\omega_2)+s\mathrm i\eta\wedge\bar\eta,
 \qquad s>0,
\]
one has
\begin{equation}\label{eq:calabi-eckmann-lower-mass}
 v_-(g)=48\pi^2sB_0.
\end{equation}
Every Hermitian $h\geq c g$ satisfies
\[
 \sup_{T\in\mathscr G_h}\|k_T\|_{L^{3/2}(\mathrm dV)}
 \leq c^{-1}(48\pi^2sB_0)^{-1/3}.
\]
No balanced, pluriclosed or LCK metric exists on $X_\tau$.
\end{proposition}
\begin{proof}
Every orbit has a unique point with $|z|=|w|=1$: the needed parameter solves
\[
 \operatorname{Re}t=-r_1/2,\qquad
 \operatorname{Im}t=(r_2-Ar_1)/(2B_0).
\]
These smooth formulas give a product trivialization by $\C\times(S^3\times S^3)$, proving freeness and properness. A branch of the logarithm sets one nonzero $z$-coordinate to one, with local holomorphic coordinates $z_{\rm other}/z_{\rm chosen}$ and $w_jz_{\rm chosen}^{-\tau}$. Thus the quotient carries the asserted compact complex structure. Logarithm branches account for the residual fiber periods.

The invariant form $\eta$ annihilates the holomorphic action generator, so it descends. On an elliptic fiber it is $\mathrm d\chi$, where $\chi=\tau\log v_1-\log v_2$ has lattice generated by $2\pi\mathrm i\tau$ and $-2\pi\mathrm i$. Hence $A_E=8\pi^2B_0$. Direct differentiation gives
\[
 \begin{aligned}
 F&=2\pi\mathrm i(\tau\omega_1-\omega_2),&
 \alpha&=-2\pi B_0\omega_1,\\
 \beta&=2\pi(A\omega_1-\omega_2),&
 Q&=-8\pi^2A\,\omega_1\wedge\omega_2>0.
 \end{aligned}
\]
Since $\int_B(\omega_1+\omega_2)^2=2$, \cref{thm:elliptic-curvature-endpoint} gives the mass, endpoint and pluriclosed obstruction.

Moreover $\mathrm d\operatorname{Re}\eta=-2\pi B_0\pi^*\omega_1$. This nonzero semipositive exact form excludes every balanced metric by the same Stokes pairing as above, hence also every K\"ahler metric. The quotient is simply connected. A hypothetical LCK metric therefore has exact Lee form $\mathrm df$; multiplying it by $\mathrm e^{-f}$ gives a K\"ahler metric, a contradiction. In particular no Vaisman metric exists.
\end{proof}

The curvature criterion controls bounded potentials and the full Gauduchon current family in these threefold geometries. Changing-sign curvature, primitive elliptic bundles and higher-dimensional analogues require further estimates.

